\clearpage
\SourceLeaf{451}
\begin{center}\large [ABWICKELUNGSFÄHIGE FLÄCHEN.]\end{center}
\begin{center}[1.]\end{center}
\noindent\hspace*{1.2em}1825 Dec. 4.\quad Das Nachdenken über die Theorie der abwicklungsfähigen Flächen hat uns folgendes elegante analytische Theorem an die Hand gegeben.

Ist $z$ eine solche Function von $x,y$, dass
\[
\frac{\partial\partial z}{\partial x^2}.\frac{\partial\partial z}{\partial y^2}
=\left(\frac{\partial\partial z}{\partial x\partial y}\right)^2,
\]
\noindent so lassen sich statt $x,y$ zwei andere veränderliche Grössen $t,u$ einführen, so dass $x,y,z$ als Functionen von $t,u$ ausgedrückt, in Beziehung auf $u$ linearische Functionen sind. In der That kann man für $u$ jede beliebige linearische Function von $x,y,z$ wählen.

\Needspace{7\baselineskip}
\begin{center}[2.]\end{center}
\begin{center}\textsc{Gauss} an \textsc{Olbers}. Göttingen, Juli 1828.\end{center}
\noindent\hspace*{1.2em}. . . . . Die Übersetzung des Artikels der Göttingischen Gelehrten Anzeigen über meine die krummen Flächen betreffende Abhandlung rührt, wie mir \textsc{Schumacher} sagte, von \textsc{Baily} her. Er sagte mir zugleich von dem Ausfall des p.\ \textsc{Fayolles} und wünschte etwas, was \textsc{Baily} zur Erwiederung brauchen könne,
